% ============================================================================
% CAPÍTULO 21: METODOLOGÍA RED TEAM — 13 RONDAS DE AUDITORÍA ADVERSARIAL
% ============================================================================
\chapter{Metodología Red Team: 300+ Ciclos de Auditoría Adversarial}


\textit{El Bulldog no ataca al dueño.\\
Ataca el código que el dueño escribió.\\
La diferencia es la única razón por la que el código mejora.} --- --- Protocolo Red Team, Constitución POLYDIM

\section{Introducción: La Auditoría Adversarial como Metodología de Investigación}

La metodología Red Team de POLYDIM invierte el paradigma habitual de desarrollo de software:
en vez de escribir código y luego buscar bugs, se \textit{buscan bugs primero} y se
escribe el código que los corrige. El proceso es:

\begin{enumerate}
\item \textbf{Atacar:} Un agente adversarial (Red Team) ataca el código existente buscando
bugs en 5 dimensiones: concurrencia, numerología, FFI, plataforma, y semántica del protocolo.

\item \textbf{Documentar:} Cada bug encontrado recibe un ciclo numerado (C1--C300+) con:
la descripción exacta del bug, el escenario de reproducción, y el parche propuesto.

\item \textbf{Corregir:} El código se corrige siguiendo el parche propuesto.

\item \textbf{Verificar:} La corrección se verifica con un test específico que reproduce
el bug y confirma que el parche lo elimina.

\item \textbf{Repetir:} El Red Team ataca el código corregido, buscando bugs en los
propios parches (Ronda 3: Ciclos 26--30 atacaron específicamente los parches de las
Rondas 1--2).
\end{enumerate}

\section{Las 13 Rondas: Estructura y Progresión}

\begin{tabular}{llll}p{8cm}}
\hline
\textbf{Ronda} & \textbf{Ciclos} & \textbf{Foco} \\
\hline
1 & P0-1 -- P2 (16 bugs) & Bugs críticos de argtypes, build, UAF básico, SLERP antipodal \\
2 & C1--C25 (25 ciclos) & Quiescencia SeqCst, dim=0 UB, max-seq reading, dtype/endian \\
3 & C26--C35 (10 ciclos) & Red Team de los propios parches; entregables integrados \\
4 & C36--C50 (15 ciclos) & Performance, reproducibilidad, ``Zero-Copy'' honestidad \\
5 & C51--C75 (25 ciclos) & Fuzzing, property tests, mutation testing \\
6 & C76--C100 & TLA+ formal spec, Loom verification \\
7 & C101--C150 & V110 FJLT (Fast Johnson-Lindenstrauss Transform) \\
8 & C151--C200 & V111 ring buffer, lending semantics \\
9 & C201--C250 & Integration test suite, CI pipeline \\
10 & C251--C275 & Numerical edge cases: subnormales, overflow, antipodal \\
11 & C276--C300 & Protocol formalization, phase definitions \\
12 & C301--C325 & V762 final hardening, ABI contract documentation \\
13 & C326+ & Convergencia: todos los tests PASS, certificación final \\
\hline


\end{tabular}

\section{Ronda 1: Los 16 Bugs Críticos P0--P2}

Los bugs de Ronda 1 fueron los más devastadores: afectaban la corrección básica del sistema.

\subsection{P0-1: argtypes Sin Declarar en ctypes}

\textbf{Descripción:} Python ctypes sin \texttt{argtypes} declarados pasa argumentos
como enteros nativos de Python. Para un puntero a \texttt{double*}, ctypes pasa
el entero del \textit{value} del objeto Python, no la dirección de memoria. El resultado
es que el kernel C++ lee desde una dirección inválida.

\textbf{Código roto:}
\begin{verbatim}
# ROTO: sin argtypes, ctypes pasa Python integers como argumento
result = cpp_lib.slerp_v109(p_ptr, q_ptr, t_val, out_ptr, D)
# El kernel lee desde la dirección INTEGER del objeto Python
# -> SIGSEGV o datos corruptos silenciosos
\end{verbatim}

\textbf{Parche P0-1:}
\begin{verbatim}
# CORRECTO: argtypes declaran el tipo exacto de cada argumento
cpp_lib.slerp_v109.argtypes = [
    ctypes.POINTER(ctypes.c_double),  # p
    ctypes.POINTER(ctypes.c_double),  # q
    ctypes.c_double,                  # t
    ctypes.POINTER(ctypes.c_double),  # out
    ctypes.c_size_t                   # D
]
cpp_lib.slerp_v109.restype = ctypes.c_int32
\end{verbatim}

\subsection{P1-4: SLERP Antipodal --- Norma 1.2247}

\textbf{Descripción:} La fórmula SLERP~\eqref{eq:slerp} es indefinida cuando
$\langle\mathbf{p}, \mathbf{q}\rangle = -1$ (vectores antipodales), porque
$\Omega = \pi$ y $\sin(\pi) = 0$, haciendo la fórmula $0/0$.

Para vectores \textit{casi} antipodales, la fórmula produce un vector con norma
$1.2247... = \sqrt{3/2}$ en lugar de $1.0$.

\textbf{Análisis matemático:}
Sean $\mathbf{p} = (1, 0, 0)$ y $\mathbf{q} = (-1+\varepsilon, 0, 0)/\|{(-1+\varepsilon,0,0)}\|$.
Para $\varepsilon \to 0^+$:
\begin{equation}
\Omega \to \pi^-, \quad \sin\Omega \to 0^+
\end{equation}
La fórmula SLERP:
\begin{equation}
\gamma(0.5) = \frac{\sin(0.5\Omega)\mathbf{p} + \sin(0.5\Omega)\mathbf{q}}{\sin\Omega}
\to \frac{\sin(\pi/2)(\mathbf{p} + \mathbf{q})}{\sin\pi}
= \frac{1 \cdot (\mathbf{p} + \mathbf{q})}{0} \to \text{división por cero}
\end{equation}

En aritmética de punto flotante, esto no produce NaN directamente sino un vector
con norma $\approx \sqrt{3/2}$ por acumulación de errores de cancelación.

\textbf{Parche P1-4 --- Gram-Schmidt para el caso antipodal:}
\begin{verbatim}
// Detectar caso antipodal: |<p,q>| > 1 - epsilon_antipodal
double dot_pq = dot_product(p, q, D);
if (dot_pq < -1.0 + 1e-9) {
    // Construir vector perpendicular a p usando Gram-Schmidt
    // Elegir e_j el vector canónico más perpendicular a p
    size_t j = argmin_abs_component(p, D);
    double e_j[D] = {0}; e_j[j] = 1.0;
    // v_perp = e_j - <e_j, p> * p (Gram-Schmidt)
    double proj = p[j]; // <e_j, p> = p[j]
    for (size_t i = 0; i < D; ++i) {
        v_perp[i] = e_j[i] - proj * p[i];
    }
    normalize(v_perp, D); // ||v_perp|| = 1
    // Ahora SLERP desde p hasta -p pasando por v_perp
    // gamma(0.5) = cos(pi/2)*p + sin(pi/2)*v_perp = v_perp
    return v_perp; // La geodésica por el polo norte
}
\end{verbatim}

\section{Ronda 2: Los 25 Ciclos (C1--C25)}

La Ronda 2 encontró bugs más sutiles en la concurrencia y el protocolo:

\subsection{C1: UAF por Ordering Incorrecto}

Ya analizado en detalle en el Capítulo~\ref{ch:seqlock}. El corazón del problema:
Acquire/Release no implica orden total, SeqCst sí.

\subsection{C2: Undefined Behavior con Dimensión Cero}

\begin{verbatim}
// Sin guard: dim=0 produce:
// - alloc de 0 bytes -> puntero nulo o comportamiento indefinido
// - divisiones por cero en cálculos de norma
// - bucle vacío con acumuladores no inicializados
pub fn pmtp_node_new(dim: usize) -> *mut PmtpNode {
    if dim == 0 {
        return std::ptr::null_mut(); // -2 en el ABI
    }
    if dim > DIM_MAX { // DIM_MAX = 1 << 27 = 128M elementos
        return std::ptr::null_mut();
    }
    // ... allocación segura ...
}
\end{verbatim}

\subsection{C22: Overread por dtype Erróneo}

Detallado en el Capítulo~\ref{ch:seqlock}: recibir \texttt{float32} cuando el
kernel espera \texttt{float64} produce una lectura de $2\times$ el buffer asignado.

\subsection{C23: El Contrato ABI}

El Ciclo 23 descubrió que sin documentación explícita, cada integrador implementa
sus propias interpretaciones de los códigos de retorno. La consecuencia: integración
incorrecta de retry logic que puede llevar a UAF por reintentar \texttt{free} después de éxito.

\section{Ronda 3: Atacar los Propios Parches}

El aspecto más valioso metodológicamente de la Ronda 3 fue descubrir bugs en los
parches de las Rondas anteriores:

\begin{description}
\item[Ciclo 26 (auto-crítica del C3):] El parche C3 (lectura por max-seq) fue roto
con múltiples writers: el seq se asigna atómicamente pero el \textit{commit} ocurre
en orden no determinista. La solución correcta (consumidor con filtro monotónico)
estaba disponible desde el Ciclo 21 pero en el lado equivocado (productor).

\item[Ciclo 27 (auto-crítica del C1):] El CAS para DESTROYING en free fue reescrito
en la Ronda 2 de manera que omite la verificación cuando el estado ya es DESTROYING.
Esto permite que dos destructores simultáneos pasen la verificación.

\item[Ciclo 28 (hueco sin parche):] Un panic en medio de una operación de escritura
deja \texttt{active\_ops $\ge$ 1} permanentemente (el nodo queda inliberable).
Ninguna ronda anterior había abordado este caso. La solución: guards RAII con Drop.
\end{description}

La lección metodológica: \textbf{los parches también necesitan ser auditados}.
La auditoría de los propios parches es la práctica más valiosa de ingeniería de software.

\section{El Ciclo 38: La Matemática Correcta del Umbral}

El Ciclo 38 es el ejemplo perfecto de la diferencia entre un umbral \textit{ad hoc}
y un umbral \textit{matemáticamente correcto}:

\textbf{El umbral roto (V109):}
\begin{equation}
\text{si } \|{\mathbf{p}\|} < 10^{-15} \text{ entonces rechazar}
\end{equation}

Este umbral rechaza vectores perfectamente válidos de norma $10^{-200}$ (que
son normalizables sin pérdida en doble precisión). El único rango que es
genuinamente no normalizable es:

\begin{equation}
\text{si } \|{\mathbf{p}\|} < \text{DBL\_MIN} = 2^{-1022} \approx 2.23 \times 10^{-308} \text{ entonces rechazar}
\end{equation}

porque $1/\|{\mathbf{p}\|}$ desborda a $\infty$ si y solo si $\|{\mathbf{p}\|} < \text{DBL\_MIN}$.

\begin{theorem}[Umbral Matemáticamente Correcto para SLERP]

La normalización de $\mathbf{p} \in \mathbb{R}^D$ en doble precisión produce overflow si y solo si
$\|{\mathbf{p}\|}_2 < \text{DBL\_MIN} = 2^{-1022}$.
Para $\|{\mathbf{p}\|}_2 \ge \text{DBL\_MIN}$, la operación
$\hat{\mathbf{p}} = \mathbf{p}/\|{\mathbf{p}\|}_2$ produce un resultado finito.
\end{theorem}

\begin{proof}
$1/\|{\mathbf{p}\|}_2$ es representable en doble precisión si y solo si
$1/\|{\mathbf{p}\|}_2 \le \text{DBL\_MAX} = (2 - 2^{-52}) \cdot 2^{1023}$.
Esto es equivalente a $\|{\mathbf{p}\|}_2 \ge 1/\text{DBL\_MAX} = \text{DBL\_MIN}$.
\end{proof}

\section{Metodología: La Pirámide de Verificación}

POLYDIM adopta una pirámide de verificación de cuatro niveles:

\begin{description}
\item[Nivel 1 --- Unit Tests (Rápido):] Cada función tiene tests parametrizados
que cubren los casos normales y los casos límite identificados por el Red Team.
Tiempo: segundos.

\item[Nivel 2 --- Property Tests (Minutos):] Las propiedades matemáticas
(invarianza de norma, monotonicidad del SEQLock, correctitud del Gram-Schmidt)
se verifican con \textit{property-based testing} (Hypothesis en Python, proptest en Rust).
Tiempo: minutos.

\item[Nivel 3 --- Loom (Exhaustivo, Horas):] Todas las entrelazados posibles de
operaciones atómicas se exploran exhaustivamente para el protocolo de quiescencia.
Tiempo: horas.

\item[Nivel 4 --- Miri (UB Detection, Horas):] El intérprete Miri de Rust detecta
Undefined Behavior que los compiladores optimizan silenciosamente. Cubre:
use-after-free, reads de memoria no inicializada, violaciones de aliasing.
\end{description}

\textbf{[CERTIFICADO]} 
\textbf{Estado V762:}
Nivel 1: 5/5 suites PASS, Exit Code 0.
Nivel 2: property tests para Rodrigues, SLERP, Neumaier --- todos PASS.
Nivel 3: Loom --- 0 violaciones en $> 10^6$ entrelazados.
Nivel 4: Miri --- pendiente (requiere Linux, integración en CI planificada para V770).

